\section*{Chapitre \ref{ch:s1:carry-across-asset-classes} --- Le carry à travers les classes d'actifs}

\begin{solution}{exo:s1:carry-across-asset-classes:1}
$\ln(100/98) \times 252/21 = 0.242$~: 24.2~\% par an, une courbe fortement en backwardation.
\end{solution}

\begin{solution}{exo:s1:carry-across-asset-classes:2}
$5~\% - 1~\% = 4~\%$ par an pour la devise~; $2~\% - 4.5~\% = -2.5~\%$ par an pour le contrat à terme sur indice~: une position longue sur contrat à terme sur actions paie le taux de financement et
reçoit les dividendes.
\end{solution}

\begin{solution}{exo:s1:carry-across-asset-classes:3}
Son cours au comptant suit un cycle annuel que la courbe anticipe~: la pente à un mois mesure surtout vers où va la saison (les cours d'hiver au-dessus
de ceux d'été), et non le carry qui serait gagné si toute la courbe restait inchangée. Sur les matières premières saisonnières synthétiques, la pente à un mois a une
corrélation de 0.05 avec le carry réel~; une pente à douze mois compare la même saison et annule le cycle.
\end{solution}

\begin{solution}{exo:s1:carry-across-asset-classes:4}
Un écart type sur cent jours vaut $10~\% \times \sqrt{100/252} = 6.3~\%$, donc la perte représente $26.3/6.3 = 4.2$ écarts types. L'espérance valait $0.95 \times
10~\% \times 100/252 = 3.8~\%$.
\end{solution}

\begin{solution}{exo:s1:carry-across-asset-classes:5}
La diversification repose sur la faible corrélation entre les classes, qui vaut en moyenne~; lors du krach, toutes les classes ont perdu en même temps, si bien que le portefeuille diversifié, ramené
à la même volatilité de 10~\%, compte tenu de sa faible corrélation moyenne, porte plus de risque brut que n'importe quelle classe seule et perd davantage.
\end{solution}

\begin{solution}{exo:s1:carry-across-asset-classes:6}
Le rendement excédentaire espéré égale le carry plus l'appréciation espérée des cours. Le carry est connu d'avance~; la question est la part que les cours lui reprennent. Si aucune,
le portefeuille de carry encaisse tout son écart de carry~; si les cours ont bougé pour l'annuler, comme le veut la parité des taux d'intérêt non couverte, il ne gagne rien. Dans l'univers synthétique
les ratios de Sharpe doublent ou triplent à peu près entre $\lambda = 0.5$ et $\lambda = 1$.
\end{solution}

\begin{solution}{exo:s1:carry-across-asset-classes:7}
Le portefeuille de matières premières passe de 0.59 à 0.34 et le portefeuille diversifié de 0.95 à 0.82~: pour les trois matières premières saisonnières, la pente à un mois classe la
saison plutôt que le carry, si bien qu'un tiers du portefeuille de matières premières négocie du bruit.
\end{solution}

\begin{solution}{exo:s1:carry-across-asset-classes:8}
Dix années calmes disent peu de chose d'une stratégie dont les pertes surviennent lors de rares krachs~: le portefeuille synthétique de devises, à 10~\% de volatilité, a perdu 25--29~\% lors de chacun
krach. Mesurer l'asymétrie et l'exposition au risque global, appliquer au portefeuille des tests de résistance fondés sur des débouclages passés, et le dimensionner en fonction du krach, non du calme.
\end{solution}

\begin{solution}{pb:s1:carry-across-asset-classes:1}
\begin{enumerate}
\item Le rendement si les cours restent inchangés, $(S - F)/(F\tau)$~; le déport à terme, le rendement du dividende moins le financement, le rendement plus le roll-down, le rendement de roll.
\item Le rendement excédentaire espéré égale le carry plus l'appréciation espérée des cours.
\item Position longue sur les actifs à carry élevé et courte sur ceux à carry faible d'une classe, classés par rang~; portefeuilles de carry de plusieurs classes à risque égal.
\item Sous forme de pente logarithmique entre les deux premiers contrats~; pour les matières premières saisonnières, la pente mesure la saison (corrélation 0.05), d'où l'on retient une
pente à douze mois.
\item Rangs centrés à exposition brute unitaire, coûts de 2, 1, 2 et 4 points de base par unité négociée, chaque classe ramenée à 10~\% dans l'échantillon~; le portefeuille diversifié pondéré également
et ramené de la même façon.
\item $\lambda = 0.5$~: 0.17, 0.22, 0.03, 0.27 et 0.34~; $\lambda = 1$~: 0.28, 0.53, 0.49, 0.59 et 0.95.
\item Quatre portefeuilles à faible corrélation moyenne se combinent en un seul, dont le ratio de Sharpe vaut environ le double de celui de la classe moyenne.
\item Ses krachs reprennent le carry qu'il gagne.
\item Une perte soudaine lorsque les actifs à carry élevé baissent ensemble~; des mouvements à asymétrie négative entre devises à taux élevé et devises à taux faible, issus du débouclage lorsque la liquidité de financement
se contracte.
\item Les devises à taux élevé sont plus exposées à un facteur de risque global, si bien que les investisseurs en carry détiennent du risque global dans les mauvais moments.
\item Devises~: 29.0~\%, 24.9~\% et 28.0~\%~; diversifié~: 12.9~\%, 26.3~\% et 13.1~\% ($\lambda = 1$).
\item Le débouclage des positions de carry sur devises en début d'août 2024 et la brève remontée du yen.
\item \textbf{Résultat nommé.} Si tout le carry est gagné, le ratio de Sharpe du portefeuille diversifié vaut 0.95, contre 0.28--0.59 pour les classes prises isolément~; lors du second
krach, il a perdu 26.3~\% à 10~\% de volatilité, davantage que chaque classe prise seule (devises 24.9~\%, actions 12.4~\%, matières premières 11.3~\%, obligations 4.1~\%).
\item Leurs rendements ne sont pas corrélés en moyenne et sont opposés lors des krachs~: le \omterm{def:s1:trend-following:trend}{suivi de tendance} gagne, le carry perd.
\item Une moyenne de 1.2~\% par an, une backwardation observée sur 51.7~\% des jours~; 24.8~\% par an pour un ratio de Sharpe de 0.61.
\item $+126~\%$ le 14 janvier 1991 et $-575~\%$ le 21 avril 2020, le lendemain du jour où le contrat de mai a clôturé sous zéro.
\item Carry calculé à partir des cours connus à l'époque, des contrats réellement détenus, des coûts et des périodes de krach dans l'échantillon.
\item Carry de la courbe du WTI.
\item Les \omterm{def:s1:carry-across-asset-classes:carry}{stratégies de portage} (carry) de toutes les classes font mal ensemble lors des récessions mondiales.
\item Pour détenir des actifs qui perdent dans les mauvais moments~: le risque global et le risque de krach.
\end{enumerate}
\end{solution}

\begin{solution}{iq:s1:carry-across-asset-classes:1}
L'écart de taux d'intérêt (déport à terme)~; le rendement excédant le taux de financement, plus le roll-down~; le rendement de roll, soit la pente de la courbe depuis le
comptant ou le contrat proche jusqu'au contrat détenu.
\end{solution}

\begin{solution}{iq:s1:carry-across-asset-classes:2}
La parité des taux d'intérêt non couverte veut que les devises à taux élevé se déprécient en moyenne de l'écart de taux, si bien que le carry est repris. Dans les faits, elles se sont
dépréciées de moins, souvent pas du tout en moyenne, avec des krachs ponctuels~; l'écart est la prime de carry.
\end{solution}

\begin{solution}{iq:s1:carry-across-asset-classes:3}
Reconnaître le début d'un débouclage~: réduire l'exposition brute jusqu'aux limites fixées à l'avance, en commençant par les paires les plus encombrées, plutôt qu'attendre que les pertes
ne l'imposent~; éviter d'alimenter le mouvement~; réévaluer lorsque la volatilité retombe.
\end{solution}

\begin{solution}{iq:s1:carry-across-asset-classes:4}
Avec des tests de résistance fondés sur des débouclages passés, appliqués aux positions du jour, l'asymétrie de ses rendements sur des échantillons plus longs et sur d'autres marchés, son bêta aux
facteurs de risque globaux, et le positionnement des autres traders de carry.
\end{solution}

\begin{solution}{iq:s1:carry-across-asset-classes:5}
Points de terme et taux d'intérêt, prévisions de dividendes et cours des contrats à terme sur indices, courbes obligataires et données sur le moins cher à livrer (CTD), courbes des matières premières par contrat,
calendriers de roll et prévisions de volatilité~; cotations périmées, contrats manquants, jours fériés, erreurs de roll et conversions de devises peuvent toutes fausser un carry.
\end{solution}

\begin{solution}{iq:s1:carry-across-asset-classes:6}
$\ln F = \ln S - c\tau$ et $d\tau = -dt$, donc $dF/F = dS/S + c\,dt$ (auquel s'ajoute un terme du second ordre qui s'annule en espérance avec $\ln$). Avec $S$ une
martingale, $E[dF/F] = c\,dt$. Avec $E[dS/S] = -(1 - \lambda)c\,dt$, $E[dF/F] = \lambda c\,dt$~: la part $\lambda$ du carry est gagnée.
\end{solution}
